\uuid{wziF}
\titre{Maximum/minimum}
\auteur{bodin ; hanani ; mzari}
\organisation{exo7}
\licence{CC BY-NC-SA 4.0}
\datecreate{2020-01-20}
\niveau{L1}
\module{Analyse}
\chapitre{Continuité, limite et étude de fonctions réelles}
\sousChapitre{Etude de fonctions}
\difficulte{2}
\theme{Maximum, bijection, maximum, minimum}
\exo7qcmid{lille-1.continuite.027}

\contenu{
\begin{question}
\qclassification{123.04}
\qtags{motcle=maximum/minimum}
Soit $f(x) = x \sin(\pi x) - \ln(x) - 1$ définie sur $]0,1]$.
Quelles sont les assertions vraies ?
\begin{answers}
    \bad{$f$ est bornée et atteint ses bornes.}

    \bad{$f$ est majorée.}

    \good{$f$ est minorée.}

    \good{Il existe $c \in ]0,1]$ tel que $f(c)=0$.}
\end{answers}
\begin{explanations}
Attention l'intervalle de définition n'est pas fermé borné. Par contre la limite de $f$ en $0$ en $+\infty$ et $f(1) = -1$. On en déduit que $f$ n'est pas majorée, par contre elle est minorée, et par la théorème des valeurs intermédiaires, elle s'annule.
\end{explanations}
\end{question}
}
